% git_hash=d53d1b9f, timestamp=2026-08-21 15:27:59 EDT
% book_springer/lectures_source/Lesson02.1_From_Data_Science_To_Decision_Science.smd

\chapter{Introduction}

% From: book_springer/lectures_source/Lesson02.1_From_Data_Science_To_Decision_Science.smd:8 '# Why Decisions, Not Predictions'
\section{Why Decisions, Not Predictions}
\label{chap:why-decisions-not-predictions}

% From: book_springer/lectures_source/Lesson02.1_From_Data_Science_To_Decision_Science.smd:24 '* From Data Science to Decision Science'
\subsection{From Data Science to Decision Science}
\label{sec:data-science-to-decision-science}

This chapter works through why a purely predictive pipeline falls short in many business
settings and contrast it with the goals, methods, and outputs decision science supplies
in its place

Each section builds on the last, moving from the four ingredients a decision system
needs, to the sharper predictive-versus-causal distinction, to the maturity ladder
that measures how many of those ingredients an organization has actually put in place

% From: book_springer/lectures_source/Lesson02.1_From_Data_Science_To_Decision_Science.smd:10 '* Why Traditional ML Falls Short'
\textbf{Why Traditional ML Falls Short} --- Traditional machine learning was built
to answer one question well: given a fixed distribution of data, what does the pattern
say will happen? \emph{Predictive accuracy} on a held-out sample is the quantity
that gets tuned and reported, and by that yardstick the field has been
remarkably successful
%
The trouble is that turning a prediction into a decision asks for more than an
accurate guess about the future, and standard ML pipelines are silent on four things
a decision requires:

\begin{itemize}
  \item \emph{Causality}: a predictive model learns correlation, not the
    mechanism that produces an outcome, so it cannot say what happens if the world
    is pushed rather than merely observed

  \item \emph{Uncertainty}: a point prediction carries no quantification of what
    the model does not know, leaving a decision-maker unable to tell a confident
    estimate from a guess

  \item \emph{The business objective}: a prediction is not a decision; something
    still has to translate a forecast or a score into an action, and that translation
    step is exactly where the business goal lives

  \item \emph{Dynamics}: the world reacts to the actions that a prediction informs,
    so the distribution a model was trained on can no longer be assumed to hold
    once the model is deployed
\end{itemize}

Each of these four gaps, on its own, is enough to turn a technically accurate
model into a poor decision-maker. A churn predictor with excellent AUC still does
not say which lever reduces churn; a demand forecast reported as a single number
hides whether the business is looking at a narrow range or a wide one; a recommender
trained on clicks optimizes clicks, not the retention or wellbeing the business
actually wants; and a fraud model trained on last year's attack patterns degrades
the moment attackers adapt to it. None of these failures shows up as an error in
the model's training loss, which is precisely what makes them dangerous: the
metrics that data science optimizes keep looking good while the decisions built
on top of them get worse. The rest of this chapter, and the three that follow it,
walk through each of these four gaps in turn, building the case for why a
decision system needs causal reasoning, calibrated uncertainty, an explicit
objective, and an awareness of its own feedback effects, not just an accurate
predictor

\textbf{From Data Science to Decision Science} --- Data science and decision
science are frequently used as if they were the same discipline wearing two
names, but they optimize for different things end to end: the question they ask,
the methods they reach for, the output they produce, the metric they report, and
even the data they need. Table~\ref{tab:data-science-vs-decision-science} lays out
the contrast dimension by dimension

\begin{table}[!t]
  \caption{Data science and decision science compared across six dimensions.}
  \label{tab:data-science-vs-decision-science}
  \begin{tabular}{p{2.3cm}p{4.3cm}p{4.7cm}}
    \hline
    \noalign{\smallskip} Dimension                        & Data Science           & Decision Science                      \\
    \noalign{\smallskip}\svhline\noalign{\smallskip} Goal & Predict outcomes       & Choose best actions                   \\
    Question                                              & What will happen?      & What should we do?                    \\
    Method                                                & Statistical models, ML & Causal models, decision theory        \\
    Output                                                & Forecast               & Action, uncertainty, expected outcome \\
    Metric                                                & Accuracy, AUC, RMSE    & Business outcome, ROI                 \\
    Data                                                  & Observations           & Observations + causal structure       \\
    \noalign{\smallskip}                                   \hline
    \noalign{\smallskip}
  \end{tabular}
\end{table}

Reading down each column in Table~\ref{tab:data-science-vs-decision-science} shows
a consistent pattern: data science stops at the forecast, while decision science
treats the forecast as raw material for an action. A data scientist reports that
accuracy improved from 91\% to 93\%; a decision scientist reports that a policy change
is expected to add \$2M in revenue with a quantified range of outcomes The two
are not competitors, since decision science needs the statistical and machine learning
tools of data science as inputs, but they answer different questions and, critically,
require different kinds of data. Predicting an outcome needs only observations of
what happened. Choosing an action needs, in addition, some model of causal
structure: what would have happened under a different action, which is a
quantity no amount of purely observational data resolves on its own

The practical stakes of this distinction are not abstract. Teams that keep
shipping predictions when the business is asking for decisions leave value on
the table, not because their models are wrong, but because a forecast alone does
not tell anyone which lever to pull, how confident to be in pulling it, or what
it costs to be wrong

This gap is also an organizational one, not just a technical one. A data science
team is typically staffed, evaluated, and incentivized around predictive metrics,
accuracy, AUC, RMSE, because those are the metrics a supervised learning pipeline
naturally produces and that a manager can compare across projects. A decision
science team has to be evaluated against business outcomes instead, ROI, cost
saved, revenue generated, which are slower to observe, harder to attribute to a single
model, and often require the causal reasoning this book builds toward just to
measure correctly. Many organizations that believe they have built a decision science
capability have, in practice, only renamed their data science team, keeping the
old predictive metrics while adopting the new vocabulary. The gap between data science
and decision science is fundamentally a gap in what gets optimized and what gets
output, and closing it is the organizing theme of this book

% From: book_springer/lectures_source/Lesson02.1_From_Data_Science_To_Decision_Science.smd:48 '* Causal vs Predictive Questions'
\textbf{Causal vs Predictive Questions} --- The distinction between data science
and decision science becomes concrete once it is applied to a specific business
question, because almost every business question has two forms: a predictive
form and a causal form, and only one of them says what to do. Table~\ref{tab:predictive-vs-causal-questions}
lines up four common examples side by side

\begin{table}[!t]
  \caption{The same business situation, phrased as a predictive question and as a
  causal question.}
  \label{tab:predictive-vs-causal-questions}
  \begin{tabular}{p{5.7cm}p{5.7cm}}
    \hline
    \noalign{\smallskip} Predictive Question                                     & Causal Question                          \\
    \noalign{\smallskip}\svhline\noalign{\smallskip} Which customers will churn? & What action prevents churn?              \\
    What will sales be next quarter?                                             & How would a 10\% price cut affect sales? \\
    Which students will fail the exam?                                           & Does tutoring improve exam scores?       \\
    Which transactions are fraudulent?                                           & What policy change reduces fraud?        \\
    \noalign{\smallskip}                                                          \hline
    \noalign{\smallskip}
  \end{tabular}
\end{table}

Each row in Table~\ref{tab:predictive-vs-causal-questions} is really one
underlying business problem asked two ways. ``Which customers will churn?''
ranks customers by risk; ``what action prevents churn?'' names an intervention
and its effect. The predictive column can be answered from observational data alone,
using the same supervised-learning toolkit that data science already supplies, and
its answers support monitoring and alerting: flag the accounts at risk, watch the
trend line, raise an alarm when a metric crosses a threshold. The causal column asks
a fundamentally different kind of question, one that requires either an
experiment or a causal model built from domain structure, because it is asking what
would happen under a change nobody may have tried yet

This is not a minor technical distinction. A model that answers only the predictive
column can be extremely accurate and still be useless for policy design, because
ranking who is at risk is not the same as knowing what to do about it

The practical reason organizations gravitate toward the predictive column first is
that it is genuinely cheaper to answer: fitting $\Pr(Y \mid X)$ from logged data
requires no experiment, no willingness to intervene on real customers or real prices,
and no causal assumptions that a stakeholder has to sign off on. The causal column
asks for exactly the ingredients a predictive pipeline was never built to supply,
either a randomized experiment, which costs time and carries its own risk, or a
causal model built from domain knowledge about how the system actually works, which
the data alone cannot validate. That asymmetry in cost is precisely why so many
organizations answer the predictive question and stop, quietly treating the ranked
list of at-risk customers as if it were already a plan of action, when in fact
it is only the first half of one. This distinction is fundamental: predictive
questions need only observational data and support monitoring, while causal questions
need causal models or experiments and support the decisions and policy design that
monitoring alone cannot provide

% From: book_springer/lectures_source/Lesson02.1_From_Data_Science_To_Decision_Science.smd:72 '* The Analytics Maturity Ladder'
\textbf{The Analytics Maturity Ladder} --- Organizations do not arrive at decision
science all at once; they climb toward it through four levels of increasing
analytical sophistication, summarized in Table~\ref{tab:analytics-maturity-ladder}

\begin{table}[!t]
  \caption{The four levels of the analytics maturity ladder, from descriptive reporting
  to full decision optimization.}
  \label{tab:analytics-maturity-ladder}
  \begin{tabular}{p{2.3cm}p{4.1cm}p{4.9cm}}
    \hline
    \noalign{\smallskip} Level                                      & Question             & Tools                   \\
    \noalign{\smallskip}\svhline\noalign{\smallskip} 1. Descriptive & What happened?       & Dashboards, reports     \\
    2. Predictive                                                   & What will happen?    & ML, forecasting         \\
    3. Causal                                                       & Why? What if we act? & DAGs, do-calculus       \\
    4. Decision                                                     & What should we do?   & Causal models + utility \\
    \noalign{\smallskip}                                             \hline
    \noalign{\smallskip}
  \end{tabular}
\end{table}

Level~1 is the dashboard: it summarizes what already happened, with no attempt to
anticipate the future. Level~2 is where most of the modern machine learning
industry lives, forecasting what will happen next from historical patterns Level~3
asks why something happened and what would happen under an intervention,
requiring the causal machinery, directed acyclic graphs and do-calculus, that a
purely predictive pipeline never needed. Level~4 closes the loop by combining a causal
model with an explicit utility function to recommend what to do, not just what
to expect

The problem this ladder exposes is that most organizations are stuck at Level~2:
they have invested heavily in forecasting and prediction infrastructure, and their
models predict well, yet nobody in the pipeline can answer what the organization
should actually do differently as a result. Part of why organizations plateau there
is momentum: the tooling, the hiring pipeline, and the internal expertise built
up around Level~2 all reinforce building yet another predictive model, since predictive
modeling is the skill the team already has, while Level~3 and Level~4 work calls
for a different set of tools, causal graphs, do-calculus, and utility functions,
that a typical machine learning curriculum does not emphasize

The jump from Level~2 to Levels~3 and~4 is not an incremental improvement in model
accuracy; it is a change in what kind of question the pipeline is built to
answer at all. Level~3 requires being explicit about the causal structure connecting
actions to outcomes, structure a purely predictive model never needs to
represent, let alone get right. Level~4 requires, in addition, an explicit
utility function that says how to weigh one outcome against another, which forces
a business to make its trade-offs precise rather than leaving them implicit in whichever
proxy metric happened to be easiest to log. Climbing to Levels~3 and~4 is the
jump this book makes, moving the reader from building prediction pipelines to building
decision pipelines, and every subsequent chapter contributes a piece of that
climb

% From: book_springer/lectures_source/Lesson02.1_From_Data_Science_To_Decision_Science.smd:96 '# The Cost of Ignoring Causality'
\section{The Cost of Ignoring Causality}
\label{chap:cost-ignoring-causality}

% From: book_springer/lectures_source/Lesson02.1_From_Data_Science_To_Decision_Science.smd:98 '## When Correlation Misleads'
\subsection{When Correlation Misleads}
\label{sec:when-correlation-misleads}

The first three sections of this chapter show that a purely predictive model,
however well it fits the data it was trained on, answers a narrower question than
the one a decision-maker needs answered. Fitting $\Pr(Y \mid X)$ tells you what
tends to co-occur with what; it does not tell you what happens if you reach in and
change $X$ yourself, and it says nothing at all about the cases history never
let you observe

% From: book_springer/lectures_source/Lesson02.1_From_Data_Science_To_Decision_Science.smd:100 '* Correlation Encodes Confounding as Causation'
\textbf{Correlation Encodes Confounding as Causation} --- A purely predictive
model learns $\Pr(Y \mid X)$ from historical data, and in doing so it absorbs every
association present in that data, whether or not the association reflects a
causal mechanism. This is not a flaw in the fitting procedure; it is what
fitting $\Pr(Y \mid X)$ means. The model answers \emph{``what will happen?''}, never
\emph{``what should we do?''}, and it implicitly assumes the world stays fixed even
though the whole point of acting on a prediction is to change the world the prediction
was made about

The canonical illustration is ice cream and drowning. Across a summer season,
ice cream sales and drowning deaths are highly correlated: as one rises, so does
the other. A predictive model built on this pattern is entirely accurate: more
ice cream sold predicts more drownings, and that association will keep holding as
long as the data-generating process stays the same. But the decision built on
top of that prediction, ban ice cream sales to cut drownings, is disastrous, because
the true driver of both quantities is hot weather. Hot weather sends more people
to buy ice cream and, independently, more people into the water, where a fraction
of them drown. Ice cream sales and drowning deaths move together only because
they share a common cause, a \emph{confounder}, not because one causes the other

This is why the observational and interventional quantities can disagree even
when the model that produced $\Pr(Y \mid X)$ is a perfect fit to the data it saw:
\begin{equation}
  \label{eq:obs-vs-int}\Pr(Y \mid X) \neq \Pr(Y \mid do(X))\;
\end{equation}
The left-hand side of Eq.~\eqref{eq:obs-vs-int} is what a predictive model estimates:
the distribution of $Y$ conditional on having observed a particular value of $X$
The right-hand side is the distribution of $Y$ under an intervention that sets $X$
to that value directly, severing whatever else in the world was previously
responsible for $X$ taking that value. Acting on $\Pr(Y \mid X)$ as if it were $\Pr
(Y \mid do(X))$ is exactly what makes correlation-based decisions backfire, and
the gap between the two is not a modeling error to be fixed with more data, but a
structural feature of observational data that only a causal model can address

% From: book_springer/lectures_source/Lesson02.1_From_Data_Science_To_Decision_Science.smd:167 '* Missing Interventions and Counterfactuals'
\textbf{Missing Interventions and Counterfactuals} --- A correlation model answers
\emph{``what is?''}, and it can be pushed, with enough data, to answer \emph{``what
is likely to be?''}. It cannot answer \emph{``what if?''} or \emph{``why?''}, no
matter how much observational data is fed into it Two classes of business
question fall squarely into this gap and stay permanently out of reach of a
purely predictive pipeline

The first is the \emph{intervention} question: \emph{``if we cut price 10\%, how
do units sold and revenue change?''} Answering this requires knowing what would
happen under a price that, in the historical data, may never have been charged, or
was charged under circumstances that differed from today's in ways the model cannot
separate out. The second is the \emph{counterfactual} question: \emph{``would
this customer have churned had we offered a discount?''} This asks about a specific
individual, under an action that was not actually taken for them, which is a strictly
harder question than the population-level intervention question, since it
requires reasoning about what would have happened to this exact customer in a
world that did not occur

The tutoring example makes the difficulty concrete. A dataset shows that
students who receive tutoring tend to score higher on exams. Does tutoring cause
the improvement, or do already-strong students disproportionately choose to sign
up for tutoring in the first place, so that the correlation reflects self-selection
rather than the effect of tutoring itself? No amount of additional observational
data on tutored and untutored students resolves this question, because the same correlation
pattern is consistent with both explanations. Distinguishing them requires
either an experiment, such as randomly assigning tutoring, or a causal model
that encodes assumptions about how students select into tutoring in the first
place, assumptions that can then be checked, argued about, and revised, unlike
the assumptions buried inside a purely correlational fit. These intervention and
counterfactual questions require a causal model, not a more accurate predictor, since
accuracy on observed data says nothing about behavior under an action the data never
contained

% From: book_springer/lectures_source/Lesson02.1_From_Data_Science_To_Decision_Science.smd:187 '* Selection Bias and the Missing Counterfactual'
\textbf{Selection Bias and the Missing Counterfactual} --- A separate and more basic
problem sits upstream of confounding: historical data records only the decision
that was actually made, never the alternative that was not. Confounding biases
the \emph{relationship} between variables that were recorded; selection bias comes
from which units got recorded at all, a problem baked into the data before any model
is fit, rather than introduced by a poor choice of model

Loan approval data is the clearest case. A bank's historical records show repayment
outcomes only for applicants who were approved; applicants who were rejected
generate no outcome at all, because they were never given a loan to repay or default
on. A credit model trained on this data learns, accurately, what predicts
repayment among the population the bank already chose to trust. It has no direct
evidence about the population the bank rejected, because that population's
counterfactual outcomes, whether they would have repaid had they been approved,
were never recorded and never can be, short of a deliberate policy change that grants
loans to a randomized sample of otherwise-rejected applicants

This creates a subtle trap for anyone evaluating the model purely by predictive
accuracy on held-out data: the held-out data suffers from exactly the same selection
process as the training data, so a model can look accurate by every standard
metric while being systematically blind to the population the business most
needs a decision about, namely the applicants currently being turned away

The trap is subtle precisely because nothing about it shows up as a data quality
problem in the usual sense: there are no missing values, no corrupted labels, no
obviously wrong entries to flag. The data is a completely accurate record of
what happened to the people the bank chose to lend to. The distortion is entirely
in what is absent, a systematic silence about an entire population, and silence of
this kind does not trigger any of the standard checks a data science team runs before
trusting a model, since those checks are built to catch errors in the data that
exists, not the absence of data that was never generated in the first place. A
model can be accurate on the population it happened to observe and still be
blind to the population the business actually needed to reason about, and no amount
of tuning on the observed population fixes a gap in what was observed to begin with

% From: book_springer/lectures_source/Lesson02.1_From_Data_Science_To_Decision_Science.smd:205 '## Structural Failure Modes'
\subsection{Structural Failure Modes}
\label{sec:structural-failure-modes}

Confounding and selection bias are not the only ways a correlation-based pipeline
goes wrong. The next four sections survey structural failure modes that can
corrupt a causal or predictive estimate even when the usual confounding checks
have been satisfied: violations of independence across units, reversals under
aggregation, illusions created by conditioning on the wrong variable, and the
discarding of mechanism in favor of whatever the training data happens to
contain

% From: book_springer/lectures_source/Lesson02.1_From_Data_Science_To_Decision_Science.smd:207 '* Interference and Spillovers Across Units'
\textbf{Interference and Spillovers Across Units} --- Standard machine learning
and standard causal inference both typically assume that each unit's outcome
depends only on that unit's own treatment, an assumption sometimes called the stable
unit treatment value assumption. Marketplaces, pricing systems, and social
networks routinely violate it, because units in these settings do not act in isolation;
they compete for the same limited resources or influence one another directly

Consider a discount shown only to group~A of customers. If group~A and group~B
compete for the same limited supply of a product, or belong to the same social network
where purchases are visible to friends, then giving group~A a discount can
cannibalize purchases that group~B would otherwise have made, or change group~B's
behavior through social influence, even though group~B never saw the discount. Once
this happens, group~B stops being a clean, untreated counterfactual for what would
have happened to group~A without the discount, because group~B's own outcome was
altered by group~A's treatment

This is not a failure of estimation technique; it is a failure of the assumption
the estimation technique relies on. The standard formula for a treatment effect,
comparing the average outcome of a treated group to the average outcome of an
untreated group, implicitly assumes that the untreated group's outcomes would have
been the same whether or not the treated group was treated. Interference breaks
that assumption directly, and it does so silently, since nothing about the
treated and untreated groups' raw outcome data reveals that the comparison group's
baseline has itself shifted

The bias this introduces does not announce itself the way a missing feature or
an obviously wrong label would; the experiment can be run correctly, the
randomization can be clean, and the estimate can still be wrong, because the
violated assumption lives outside the data the experiment collected, in the market
or network structure connecting the groups to each other. When one unit's action
changes another unit's outcome, the usual estimate of a treatment effect,
computed by comparing a treated group to an untreated group, is biased, no
matter how accurate the underlying predictive model of each individual outcome
is. Interference has to be modeled explicitly, through network structure or
market-clearing constraints, rather than assumed away

% From: book_springer/lectures_source/Lesson02.1_From_Data_Science_To_Decision_Science.smd:224 '* Simpson's Paradox'
\textbf{Simpson's Paradox}

\begin{definition}
  \emph{Simpson's paradox} occurs when a trend present within every subgroup of
  a population reverses once the subgroups are pooled together into a single
  aggregate
\end{definition}

A drug-recovery study illustrates the paradox concretely. Table~\ref{tab:simpsons-paradox}
reports recovery rates for male and female patients separately, and for the two groups
combined

\begin{table}[!t]
  \caption{A drug that helps every subgroup can appear to hurt in the pooled
  aggregate: Simpson's paradox in a recovery-rate study.}
  \label{tab:simpsons-paradox}
  \begin{tabular}{p{2.4cm}p{4.6cm}p{4.6cm}}
    \hline
    \noalign{\smallskip} Group                            & Recovery w/ drug & Recovery w/o drug \\
    \noalign{\smallskip}\svhline\noalign{\smallskip} Male & 81 / 87 = 93\%   & 234 / 270 = 87\%  \\
    Female                                                & 192 / 263 = 73\% & 55 / 80 = 69\%    \\
    Combined                                              & 273 / 350 = 78\% & 289 / 350 = 83\%  \\
    \noalign{\smallskip}                                   \hline
    \noalign{\smallskip}
  \end{tabular}
\end{table}

Reading Table~\ref{tab:simpsons-paradox} row by row, the drug helps male
patients, 93\% versus 87\%, and it helps female patients, 73\% versus 69\%. Yet
the combined row shows the drug apparently hurting, 78\% versus 83\%. The
reversal happens because gender confounds both treatment assignment and the
outcome: female patients, who recover less easily than male patients regardless
of the drug, receive the drug far more often, 263 of the 350 treated patients,
versus only 87 of 350 for male patients, so pooling stacks a mostly-female
drug group against a mostly-male no-drug group, and the harder-to-treat gender
dominates the ``drug'' column and reverses the sign of an effect that helps
everyone

The paradox is easy to dismiss as a curiosity, yet the pooled row in Table~\ref{tab:simpsons-paradox}
is exactly the number a Level~1 dashboard would surface: an overall recovery
rate, computed without regard to any subgroup structure, because that is the
number that is cheapest to compute and easiest to put in a single-line summary. Nothing
about the computation is wrong; the aggregate recovery rate really is what it
says it is. What is wrong is treating that single number as though it settled the
causal question of whether the drug helps, when the answer depends entirely on
which stratification the analyst chose to stop at. The same mechanics generalize
well past a two-category example: any variable that correlates with both
treatment assignment and outcome can drive a reversal, and the more subgroups a population
can be split along, the more directions a naive aggregate can be pulled in. A model
can look correct, or even look harmful, in aggregate and be exactly wrong about every
segment it is actually applied to. This is why a causal analysis needs to
identify the right level of stratification, guided by the causal structure connecting
treatment, subgroup, and outcome, before trusting an aggregate number

% From: book_springer/lectures_source/Lesson02.1_From_Data_Science_To_Decision_Science.smd:296 '* Berkson's Paradox'
\textbf{Berkson's Paradox}

\begin{definition}
  A \emph{collider} is a variable that is influenced by two or more independent
  causes. Conditioning on a collider, whether by filtering a dataset, splitting a
  sample, or including the collider as a covariate, opens a spurious statistical
  association between its causes that did not exist before the conditioning This
  mechanism is known as \emph{Berkson's paradox}
\end{definition}

Consider two traits, statistics skill and flattery, that both independently
influence whether someone gets promoted, with no direct relationship between the
two traits themselves. Among the full population of employees, statistics skill and
flattery are statistically independent, exactly as assumed. But look only among the
employees who were promoted, that is, condition on the collider ``promoted'',
and a negative association appears between the two traits: among the promoted,
being bad at statistics implies the person must have been a good flatterer to
make up for it, because being weak on both traits would have kept them out of the
promoted group in the first place

Nothing about statistics skill and flattery changed; what changed is the population
being examined. Filtering to the promoted subgroup, which is exactly what happens
whenever an analysis restricts itself to survivors, admitted applicants, or any outcome-selected
sample, silently opens a spurious channel of association between the two
independent causes of that outcome

Collider bias is easy to introduce by accident, because the instinct in most
predictive modeling workflows is to add any available variable as a covariate on
the assumption that more information can only help. If that variable happens to be
a collider of two other variables already in the model, adding it does not add
information; it manufactures a spurious association between them that the model then
has no way to distinguish from a genuine one. Recognizing a collider requires
knowing, or hypothesizing, the causal graph connecting the variables under study;
it is not something that can be read off the correlational structure of the data
alone, since the spurious association a collider creates looks, statistically,
exactly like a real one. Conditioning on the wrong variable, not merely omitting
the right one, is by itself enough to manufacture a correlation with no causal meaning
at all. This is why deciding what to control for in an analysis requires knowing
the causal structure, not just adding available covariates to a regression

% From: book_springer/lectures_source/Lesson02.1_From_Data_Science_To_Decision_Science.smd:316 '* Discarding Domain Knowledge and Mechanisms'
\textbf{Discarding Domain Knowledge and Mechanisms} --- A purely correlation-based
pipeline discards known structure, whether physics, economics, or an explicit
causal graph, in favor of whatever pattern happens to be present in the training
data. Standard supervised learning offers no principled way to inject such
structure as a constraint or a prior; the model is free to fit any pattern that reduces
training loss, whether or not that pattern reflects the actual mechanism
generating the outcome

The wolf-versus-husky classifier is the textbook case of this failure. A model
trained to distinguish wolves from huskies can reach high test accuracy while, on
inspection, keying almost entirely on the presence of snow in the background
rather than on any feature of the animal itself, because wolf photographs in the
training set happened to be taken in snowy settings far more often than husky photographs
were. The classifier is accurate on the test distribution for the wrong reason, and
it becomes worthless the moment conditions change: a husky photographed in snow,
or a wolf photographed without it, breaks the pattern the model actually learned

This is not an isolated anecdote; it is the generic failure mode of any model
that fits pattern without mechanism. The deeper issue is that a purely data-driven
pipeline has no slot in which to place domain knowledge even when that knowledge
is available and well established: a physicist's conservation law, an economist's
budget constraint, or a known causal graph could all, in principle, rule out
spurious patterns like the snow-background shortcut before the model ever has
the chance to learn them, but standard supervised learning offers no mechanism for
injecting such a constraint as anything other than another feature to be weighed
against all the others. A model that ignores the mechanism generating its data
inherits every artifact of how that training data happened to be collected, and no
amount of additional training data of the same kind fixes the problem, since
more data of the same kind only reinforces the spurious pattern rather than
correcting it

% From: book_springer/lectures_source/Lesson02.1_From_Data_Science_To_Decision_Science.smd:376 '# The Cost of Ignoring Uncertainty'
\section{The Cost of Ignoring Uncertainty}
\label{chap:cost-ignoring-uncertainty}

% From: '## Point Estimates in a Small-Data World'
\subsection{Point Estimates in a Small-Data World}
\label{sec:point-estimates-small-data}

The four sections below trace what goes wrong when a pipeline reports a single number
without saying how much to trust it: the spread hidden behind a point forecast, the
two different sources of uncertainty that get flattened into one, the absence of
any mechanism for a model to admit it does not know, and the statistical traps that
let unwarranted confidence accumulate silently over a project

% From: book_springer/lectures_source/Lesson02.1_From_Data_Science_To_Decision_Science.smd:378 '* Point Estimates Without Error Bars'
\textbf{Point Estimates Without Error Bars} --- Most business problems are, in
the relevant sense, small-data problems: the number of comparable historical
episodes, a product launch, a pricing change, a new market, is small relative to
how much variation the outcome can show, even when the raw row count in a database
looks large. Standard machine learning pipelines are not built with this in mind
They are built to emit a single number, and that number is typically reported
and consumed as though it were certain, with no error bars, no posterior
distribution, and no sense of how far off the estimate could plausibly be

The gap this creates is easiest to see in a forecasting example. A demand forecast
of ``1,200 units next month'' is a single number, but it hides an enormous
amount of information about how much to trust it. If the true range is 1,100 to 1,300
units, the business can order close to 1,200 units with confidence. If the true range
is 400 to 3,000 units, ordering 1,200 units is a bet, not a plan, and the
inventory decision that makes sense under the narrow range can be badly wrong
under the wide one. The forecast number itself is identical in both cases; only
the spread around it tells the decision-maker which situation they are actually
in

This is not a minor omission that a more careful analyst can patch after the fact,
because the spread is not implied by the point estimate; it has to be estimated and
reported alongside it, using tools, confidence intervals, Bayesian posteriors, ensembles,
that a purely point-estimate pipeline was never asked to produce

The same point number can imply opposite actions once its uncertainty is made
explicit. Take a prediction that a stock's price will rise 1\% If that estimate
carries a 95\% confidence interval of $1\% \pm 2\%$, the true return could easily be
negative, and the sensible response is to ignore the prediction. If the same point
estimate instead carries an interval of $1\% \pm 0.01\%$, the outcome is all but
certain, and the sensible response is to buy. The point forecast is identical,
1\%, in both cases; only the width of the interval around it tells the
decision-maker which of the two opposite actions is the right one

The small-data character of most business problems makes this worse, not better,
because the usual intuition that more data shrinks uncertainty automatically applies
most cleanly in the large-sample regime, and a great many consequential business
decisions, a product launch, a new market entry, a pricing change, do not have
that luxury. A model can still report a single confident-looking number in these
settings, since nothing about a standard training pipeline forces it to
acknowledge how thin the evidence actually is, which means the point estimate can
look identical whether it rests on a thousand comparable historical episodes or on
three. A decision-maker needs the spread of an estimate at least as much as its
center, since the center alone cannot distinguish a safe bet from a gamble

% From: book_springer/lectures_source/Lesson02.1_From_Data_Science_To_Decision_Science.smd:398 '* Epistemic vs Aleatoric Uncertainty Conflated'
\textbf{Epistemic vs Aleatoric Uncertainty Conflated} --- Even when a model does
report some notion of uncertainty, standard pipelines typically collapse two
genuinely different kinds of uncertainty into a single number \emph{Aleatoric} uncertainty
is irreducible randomness in the outcome itself: even with a perfect model and
infinite data, the outcome would still vary, because the underlying process is
stochastic, as with the outcome of a fair coin flip. \emph{Epistemic} uncertainty,
by contrast, comes from not having seen enough data, or from using the wrong model;
it reflects what the model does not yet know, rather than randomness inherent to the
world, as with a rare disease for which only ten cases have ever been recorded

The distinction matters because the two kinds of uncertainty call for different responses,
and without separating them a model cannot flag the situation where it is most
likely to be dangerously wrong: exactly off its training distribution, on inputs
unlike anything it has seen before, which is also exactly where high-stakes decisions
tend to bite hardest. A fraud-detection model that reports low predicted risk
for a novel transaction pattern might be reporting a genuinely low aleatoric
risk, or it might simply have never seen this pattern before and be reporting a confident-sounding
number that reflects epistemic ignorance rather than genuine safety. The two
cases look identical in the output but call for opposite responses: one calls for
approving the transaction, the other for routing it to a human reviewer or
collecting more data before trusting the model's answer

The two kinds of uncertainty also imply different remedies at the system level,
not just at the level of a single prediction. Aleatoric uncertainty is a property
of the outcome itself, so no amount of additional modeling effort or additional
data collection reduces it; a coin flip stays a coin flip. Epistemic uncertainty
is a property of what the model has been shown, so it responds to exactly the interventions
a decision-maker can actually take: collecting more examples from the sparse
region of the input space, gathering an additional feature that resolves the ambiguity,
or deploying a more flexible model class. Conflating the two into a single reported
number obscures which of these responses, if any, is worth pursuing for a given prediction
Only epistemic uncertainty shrinks as more data is collected, since aleatoric uncertainty
is a property of the outcome, not of the model's knowledge; treating both kinds of
uncertainty as a single number misleads not just the current decision but also the
decision of when it is worth collecting more data at all

% From: book_springer/lectures_source/Lesson02.1_From_Data_Science_To_Decision_Science.smd:421 '* No Abstention: Systems That Never Say "I Don't Know"'
\textbf{No Abstention: Systems That Never Say ``I Don't Know''} --- A standard classifier
or regressor is built to always return an answer. Given any input, including an input
that falls far outside the range of anything the model was trained on, the model
still produces a prediction, because nothing in its architecture or training procedure
includes a mechanism for declining to answer. There is no built-in notion of
abstention or deferral that would route an unusually uncertain case to a human
decision-maker instead of forcing the model to guess

A loan-approval model illustrates the cost of this gap concretely. When the
model is handed an applicant profile that looks unlike anything in its training
data, perhaps an unusual combination of income source, credit history, and requested
loan amount, it does not, and structurally cannot, respond by saying the case is
out of distribution and should be reviewed manually. It still returns a confident-looking
approval or denial, computed from whatever the model's decision boundary happens
to say about that region of the input space, a region the model has effectively no
evidence about

This is a design gap, not a training failure: no amount of additional training
on the existing distribution adds an abstention mechanism, because abstention has
to be built into the system as an explicit option, typically by comparing a case's
estimated uncertainty against a threshold and routing high-uncertainty cases
elsewhere

Building that option in is not free, since it depends on the model having a
trustworthy estimate of its own uncertainty in the first place, exactly the epistemic-versus-aleatoric
distinction from the previous section. A model with no notion of epistemic
uncertainty has nothing reliable to threshold on, so adding an abstention rule on
top of it can end up deferring the wrong cases, the ones that happen to sit near
a decision boundary for aleatoric reasons, while confidently answering the out-of-distribution
cases it should have deferred instead. Abstention and calibrated uncertainty are
therefore not two independent features to add to a pipeline; the second is a prerequisite
for the first to work as intended. A decision system that cannot say \emph{``I
don't know''} cannot be trusted with precisely the cases that matter most, since
those are disproportionately the cases where the training data gave it the least
evidence to go on

% From: book_springer/lectures_source/Lesson02.1_From_Data_Science_To_Decision_Science.smd:435 '* Overfitting and the Statistical Significance Traps'
\textbf{Overfitting and the Statistical Significance Traps} --- Without a prior
or some other constraint on how closely a model is allowed to fit its training
data, nothing bounds how much a model can overfit, memorizing noise in the
training sample rather than learning the pattern that generalizes. This basic overfitting
risk is compounded, in practice, by a set of statistical significance traps that
let unwarranted confidence build up gradually over the course of a project rather
than appearing as a single identifiable error

\begin{itemize}
  \item \emph{Peeking}: checking a test's result repeatedly over time and stopping
    the moment $p < 0.05$ is reached inflates the true false-positive rate well past
    the nominal $5\%$, because the test was effectively given many chances to
    cross the threshold by chance, not just one

  \item \emph{Multiple comparisons}: testing many metrics or many variants at once
    all but guarantees that at least one of them shows a spurious win somewhere,
    purely from the number of chances taken, even if none of the underlying effects
    are real

  \item \emph{Burning the test set}: reusing the same held-out data across many rounds
    of model tuning turns that data into training data in disguise, since each
    round of tuning is implicitly fit to whatever patterns, real or spurious,
    that specific held-out sample happens to contain
\end{itemize}

Each of these traps is individually well known, yet they compound across the
life of a typical project: a team peeks at an experiment's results a few times before
deciding to ship, checks several metrics rather than one pre-registered metric, and
reuses the same validation set across dozens of tuning iterations, with no
single step looking like a mistake in isolation

What makes the three traps dangerous together, rather than merely additive, is that
each one on its own can look like ordinary, disciplined data science practice. Checking
an experiment's dashboard periodically is standard operating procedure; tracking
several metrics rather than a single one is usually recommended, not discouraged;
and reusing a validation set across iterations is simply what iterative model
development looks like day to day. None of these steps triggers an alarm, and
none of them is a mistake in the sense of violating an explicit rule, which is exactly
why the inflated false-positive rate they jointly produce tends to surface only much
later, as a model that looked strong in development quietly underperforms once
it meets data it was not, in effect, evaluated against even once. Overfitting is
not one mistake made in one run; it accumulates from every repeated look taken
at the same data, which is why disciplined experimental protocol, not just a
better model, is what keeps it in check

% From: book_springer/lectures_source/Lesson02.1_From_Data_Science_To_Decision_Science.smd:452 '# The Cost of Ignoring the Business Objective'
\section{The Cost of Ignoring the Business Objective}
\label{chap:cost-ignoring-business-objective}

% From: book_springer/lectures_source/Lesson02.1_From_Data_Science_To_Decision_Science.smd:454 '## Proxies, Costs, and Trade-offs'
\subsection{Proxies, Costs, and Trade-offs}
\label{sec:proxies-costs-tradeoffs}

The next three sections examine the objective function itself: what happens when
it is a proxy rather than the true business goal, what happens when it treats
every error as equally costly, and what happens when a real decision has to trade
off several goals at once rather than optimize a single scalar

% From: book_springer/lectures_source/Lesson02.1_From_Data_Science_To_Decision_Science.smd:456 '* Optimizing the Wrong Objective'
\textbf{Optimizing the Wrong Objective} --- Machine learning models chase whatever
objective they are given, and that objective is frequently a proxy metric, clicks,
engagement time, predictive accuracy, rather than the underlying business goal
the proxy was chosen to stand in for, such as revenue, customer retention, or user
wellbeing. The proxy is usually chosen because it is easy to measure and easy to
optimize, while the true goal is harder to quantify or only observable with a long
delay

\begin{definition}
  \emph{Goodhart's law} states that once a proxy measure becomes a target that
  is directly optimized, it stops reliably measuring the underlying quantity it was
  originally meant to track
\end{definition}

A recommender system optimized purely for click-through rate is the standard illustration
Click-through rate is easy to log and easy to optimize directly, and a system trained
to maximize it will find the content most likely to earn a click, which is often
sensational, outrage-inducing, or otherwise attention-grabbing content. Click-through
rate rises exactly as intended, while the deeper goal the metric was supposed to
proxy for, sustained long-term engagement, or user trust in the platform, quietly
erodes, because the content that wins clicks in the short run is not the content
that keeps users satisfied over time. The model has not failed by its own lights;
it has succeeded perfectly at the objective it was actually given

The reason a proxy is chosen in the first place is usually a reasonable one:
revenue, retention, and wellbeing are often slow to observe, expensive to measure
directly, or only available with a lag long enough to make them impractical as a
training signal for a system that needs to update frequently. Clicks, by
contrast, are available instantly and in volume. The failure is not in choosing
a proxy, which is frequently unavoidable, but in treating the proxy as though it
were the goal once it has been chosen, rather than continuing to check it against
the goal it was meant to stand in for. The objective a model is trained on
becomes the objective the business actually gets, regardless of whether that was
the intent behind choosing the proxy in the first place, which makes the choice
of training objective one of the highest-leverage decisions in building a model,
not a minor implementation detail to settle after the modeling is done

% From: book_springer/lectures_source/Lesson02.1_From_Data_Science_To_Decision_Science.smd:473 '* Symmetric Losses for Asymmetric Errors'
\textbf{Symmetric Losses for Asymmetric Errors} --- Standard loss functions, log-loss,
mean squared error, treat every unit of error identically: a false positive costs
the same as a false negative of the same magnitude, and an overestimate costs
the same as an equally sized underestimate Real business errors are almost never
symmetric in this way. A fraud false negative, letting a fraudulent transaction
through, and a fraud false positive, blocking a legitimate customer's purchase, impose
very different costs on the business, and treating them as interchangeable in
the loss function bakes a mismatch into the model from the start. Stock return
prediction shows the same asymmetry from a different angle: a loss on a position
is bounded, since it cannot fall below losing the entire investment, but a run of
gains is not bounded in the same way, so a loss function that scores an
overestimate and an underestimate of the same size identically misrepresents a
risk-reward profile that is asymmetric from the start

A fraud model that reaches 95\% accuracy sounds like a success by the standard
metric, but if that accuracy comes from a model that blocks half of all legitimate
transactions in order to catch a small number of fraudulent ones, it is an
expensive and damaging decision system in practice, no matter how good its accuracy
number looks on a dashboard. The accuracy metric does not know, and cannot know on
its own, that a blocked legitimate customer and a missed fraudulent charge are
not the same size of mistake to the business

The deeper issue is that a symmetric loss is a modeling convenience, not a
business fact: it is chosen because it is differentiable, well studied, and easy
to optimize with standard tools, not because the business genuinely experiences
every error as equally costly. Recovering an asymmetric-cost model from a symmetric-loss
training run is possible only by adjusting the decision threshold after the fact,
which helps but is a strictly weaker fix than training against the true costs
directly, since a threshold adjustment cannot correct for a model that learned the
wrong decision boundary in the first place because it was never shown that the two
error types differ in cost during training. Cost asymmetry has to be built
directly into the training objective, through a weighted loss, a cost matrix, or
an explicit decision threshold chosen to reflect the relative costs, rather than
bolted on after deployment as a post-hoc adjustment to a model that was never
asked to care about the difference in the first place

% From: book_springer/lectures_source/Lesson02.1_From_Data_Science_To_Decision_Science.smd:495 '* Multi-Objective Decisions Under Constraints'
\textbf{Multi-Objective Decisions Under Constraints} --- Real decisions rarely
optimize one thing at a time; they trade off several competing goals simultaneously,
and they do so under hard constraints of budget, capacity, or legal requirement
A single scalar loss function, no matter how carefully weighted, cannot fully express
a constrained, multi-stakeholder trade-off, because collapsing several goals into
one number requires deciding, in advance and often implicitly, exactly how much of
one goal is worth sacrificing for how much of another

A hospital choosing a cancer therapy for a patient illustrates the difficulty
The choice weighs five-year survival probability, quality of life during treatment,
and financial cost simultaneously, and these three considerations do not reduce to
a single ``best'' therapy independent of how they are weighed against one another:
a therapy that maximizes expected survival may impose a much harder course of
treatment, and a cheaper therapy may trade off some survival probability for a substantially
better quality of life. Different patients, and different values placed on each dimension,
lead to different reasonable choices from the identical clinical facts

The usual engineering response to a multi-objective problem is to scalarize it,
picking fixed weights and summing the objectives into one number a standard
optimizer can maximize. Scalarizing is not wrong in principle, but it quietly
forces a decision, the relative weight placed on survival versus quality of life
versus cost, into a modeling choice made once, often by whoever built the pipeline,
rather than a choice made explicitly by the patient, clinician, or stakeholder
who actually bears the consequences of getting the weighting wrong. A genuinely multi-objective
framing keeps the trade-off visible instead, for instance by presenting a
frontier of options that are each best under some weighting, and letting the decision-maker
choose among them with the trade-off in full view rather than pre-resolved inside
a loss function nobody examined. Collapsing a genuinely multi-objective decision
into a single loss function does not solve the trade-off; it hides it, silently
baking in someone's judgment about the relative weights instead of surfacing
that judgment where it can be examined, debated, and adjusted by the people
accountable for the decision

% From: book_springer/lectures_source/Lesson02.1_From_Data_Science_To_Decision_Science.smd:511 '## From Scores to Actionable Decisions'
\subsection{From Scores to Actionable Decisions}
\label{sec:scores-to-actionable-decisions}

Even a well-specified objective is not enough if the pipeline's final output is
a bare score with no attached lever, and even a well-designed individual model is
not enough if it sits inside a pipeline optimized piece by piece rather than as a
whole. The next two sections take up these two remaining gaps between a
technically sound model and a usable decision

% From: book_springer/lectures_source/Lesson02.1_From_Data_Science_To_Decision_Science.smd:513 '* Black-Box Scores Are Not Recommendations'
\textbf{Black-Box Scores Are Not Recommendations} --- A black-box score ranks
units against one another but names no lever and gives no reason \emph{why} a given
unit received the score it did. A churn score of 0.8 tells a business that a
customer is likely to leave; it does not say what action, if any, would change that
outcome. As credit, hiring, and healthcare decisions increasingly require
explanations to satisfy regulators and justify outcomes to the people affected
by them, a bare score is decreasingly sufficient even where it once might have passed
as adequate

The difference between a score and a recommendation is the difference between ``this
customer has churn score 0.8'' and ``offering a discount lowers this customer's
churn probability by 12 points.'' The second statement names the lever, the
discount, quantifies its effect, the 12-point reduction, and identifies the unit
it acts on, this customer. It is actionable in a way the bare score is not,
because it tells the decision-maker what to do, not merely who to worry about

Notice that the second statement is a causal claim, not a predictive one: it
says what happens under an intervention, offering the discount, not merely what
tends to co-occur with a high or low churn score. This is why the gap between a score
and a recommendation cannot be closed by a better predictive model alone,
however well calibrated; closing it requires the same causal machinery introduced
earlier in this book, since naming a lever and quantifying its effect is exactly
the do-calculus question $\Pr(Y \mid do(X))$ that a purely observational score
never had to answer. A score built from correlation can rank customers correctly
while still saying nothing true about what any specific intervention would do to
any specific customer's outcome. A genuine recommendation needs both an
actionable lever and an explanation of why that lever works, neither of which a bare
predictive score provides on its own. Moving from scoring to recommending
requires causal reasoning about the effect of specific interventions, not just a
ranking of who is most at risk

% From: book_springer/lectures_source/Lesson02.1_From_Data_Science_To_Decision_Science.smd:534 '* The Cost of Solving Sub-Problems Separately'
\textbf{The Cost of Solving Sub-Problems Separately} --- A decision pipeline is frequently
built as a chain of smaller pieces, each optimized on its own: a prediction model
trained to minimize forecast error, followed by a separate, often hand-written rule
that converts the prediction into an action. Each piece can be individually well
engineered and still add up to a pipeline that performs worse than one trained
to optimize the final decision directly, an approach known as \emph{decision-focused
learning}, which has been shown to outperform this kind of piecewise construction

An inventory system illustrates the gap. A forecasting model can predict demand accurately,
minimizing forecast error on held-out data, and then hand that forecast to a fixed
reorder rule that converts the number into a purchase order The resulting orders
can still be worse, in terms of actual inventory cost, than the orders from a system
trained end to end to minimize inventory cost directly, because the forecasting
model was never asked to care about the asymmetric costs of over- and under-ordering,
only about minimizing an average forecast error that treats every unit of error
the same

The reason the split hurts is structural, not incidental: gradient information about
the true downstream cost, how much an order-quantity error actually costs the
business, never flows back into the forecasting model's training, because the
forecasting model is optimized entirely on its own, upstream objective before the
reorder rule ever sees its output. Decision-focused learning closes exactly this
gap by differentiating through the reorder rule itself, so the forecasting model
is trained with the downstream inventory cost as its actual objective, rather
than a forecast-error proxy that happens to be convenient to optimize in
isolation. Optimizing each piece of a pipeline in isolation does not add up to optimizing
the decision the pipeline as a whole is meant to produce, since the seam between
the pieces, here the fixed reorder rule, is exactly where the objective the
business cares about can leak out of the optimization entirely

% From: book_springer/lectures_source/Lesson02.1_From_Data_Science_To_Decision_Science.smd:552 '# The Cost of Ignoring Dynamics and Feedback'
\section{The Cost of Ignoring Dynamics and Feedback}
\label{chap:cost-ignoring-dynamics-feedback}

% From: book_springer/lectures_source/Lesson02.1_From_Data_Science_To_Decision_Science.smd:554 '## A World That Reacts'
\subsection{A World That Reacts}
\label{sec:world-that-reacts}

The first three sections below take up what happens once a deployed model starts
influencing the very world it was trained to describe: the distribution it was
fit to can shift, the shift can compound into a feedback loop, and in the
sharpest case the prediction itself becomes the cause of the outcome it was
meant to forecast

% From: book_springer/lectures_source/Lesson02.1_From_Data_Science_To_Decision_Science.smd:556 '* The Static-World Assumption Breaks'
\textbf{The Static-World Assumption Breaks} --- Every decision changes the world
in some way, and that change breaks the assumption, implicit in almost all
standard machine learning, that the distribution a model was trained on will still
hold once the model is deployed and acting on that world. The world a model was
fit to is not the world the model then operates in, because the model's own outputs
are part of what shapes the world going forward

This break happens in three distinct ways. \emph{Concept shift} occurs when the
relationship $\Pr(Y \mid X)$ between inputs and outcomes itself changes, so that
the same input now produces a different outcome than it used to. \emph{Covariate
shift} occurs when the input distribution $\Pr(X)$ changes while the underlying relationship
between inputs and outcomes stays the same, so the model is now seeing a
different mix of cases than it was trained on. \emph{Label shift} occurs when
the outcome distribution $\Pr(Y)$ itself changes, independent of any change in how
inputs map to outcomes

A card-fraud model trained on last year's transactions can suffer all three kinds
of shift at once, each traceable to a different cause. Attackers adopting new
tactics is a concept shift: a transaction pattern that used to be safe now
signals fraud, changing $\Pr(Y \mid X)$ itself. Contactless payments becoming far
more common is a covariate shift: the input distribution $\Pr(X)$ moves even
though fraud tactics on that channel are unchanged. A data breach that leaks
card numbers and pushes up the overall fraud rate is a label shift: the outcome
distribution $\Pr(Y)$ rises on its own, independent of any change in how inputs
map to outcomes. The model has no built-in way to detect any of the three on its
own, since it only knows the patterns it was shown during training and has no
signal that the world it is now scoring has moved on from that training period

A house-price model trained on pre-pandemic listings shows the same three shifts
playing out in a completely different market. Remote work permanently changes
what commute distance is worth to a buyer, a concept shift in how a feature maps
to price, while a move to lower interest rates raises average sale prices across
the board, a label shift in the outcome itself, with neither shift announcing
itself as an error in the model that was fit before either change occurred

The three kinds of shift are worth keeping separate because they call for different
monitoring signals: concept shift shows up as a change in the relationship
between features and outcomes even when the feature distribution looks stable, covariate
shift shows up as a change in the feature distribution even when the model's accuracy
on familiar inputs looks unchanged, and label shift shows up as a change in the
base rate of the outcome itself. A monitoring system that only tracks overall accuracy
can miss all three kinds of shift for a while, since a model can compensate for
one kind of drift by accident while a different kind is silently accumulating
underneath it. A model that encodes mechanism, the actual causal process
generating an outcome, rather than just surface pattern, degrades more
gracefully under all three kinds of shift, because mechanisms tend to be far more
stable across a distribution shift than the correlational patterns a purely predictive
model relies on

% From: book_springer/lectures_source/Lesson02.1_From_Data_Science_To_Decision_Science.smd:586 '* Feedback Loops and Bias Amplification'
\textbf{Feedback Loops and Bias Amplification} --- Predictions do not just describe
the world; once deployed, they act on it, and the world they act on then
reshapes the very data the model is later retrained on This creates a feedback
loop: the model's past decisions become part of the future training set, and
whatever bias was present in the original model gets reinforced rather than
corrected by the retraining process

A content recommender illustrates the loop directly. A recommender that surfaces
sensational content earns more clicks on that content, precisely because it is sensational;
those extra clicks then become part of the training data used to retrain the
model, so the next version of the model, now trained on data containing even
more clicks on sensational content than before, recommends more of it still
Each cycle of the loop pushes the system further in the same direction, without any
single retraining step looking like an error on its own

A second illustration comes from policing: over-policing a particular neighborhood
produces more recorded arrests there, simply because more police presence
generates more arrests for a given underlying crime rate; that arrest data then trains
a model that flags the neighborhood as higher risk, and that risk flag is used
to justify allocating even more policing to the same neighborhood, closing the loop

Both examples share a structural feature that makes feedback loops especially
hard to catch in ordinary model monitoring: every individual retraining cycle can
look like the model is doing exactly what it is supposed to do, fitting the
newest available data more closely than the previous version did. The problem is
invisible at the level of a single retraining step and only becomes visible by comparing
the trajectory of the system across many cycles, which is a comparison most
monitoring dashboards, built around a single model's current-period metrics, are
not set up to make. A feedback loop turns a small initial bias into a growing one,
purely as a consequence of the model being retrained on the outcomes of its own
past decisions, which means the bias compounds even if every individual retraining
step is executed correctly by the standards of ordinary supervised learning

% From: book_springer/lectures_source/Lesson02.1_From_Data_Science_To_Decision_Science.smd:607 '* Performativity: When the Prediction Changes the Outcome'
\textbf{Performativity: When the Prediction Changes the Outcome}

\begin{definition}
  \emph{Performativity} describes a prediction that changes the very quantity it
  is trying to predict. It is a sharper, explicitly named instance of the more
  general feedback-loop phenomenon: rather than merely influencing future training
  data indirectly, a performative prediction alters the outcome directly, as a consequence
  of being published or acted upon
\end{definition}

A published credit-risk score is a clear example. The score itself moves
borrower behavior: a borrower who learns their score is low may take on less new
credit, or may take on more in an attempt to prove creditworthiness, and lenders
who see the score adjust the terms they offer. These behavioral responses to the
score then move the true default rate that the score was originally trying to
predict, so the score's target is not a fixed quantity waiting to be estimated more
accurately; it is a quantity that shifts in response to the act of estimating
and publishing it

A navigation app makes the same effect visible in real time. Predicting a
particular route as fastest routes many drivers onto it at once, and that surge
of traffic congests the route until it is no longer the fastest option,
invalidating the very prediction that attracted the drivers to it in the first
place

Performativity sharpens the feedback-loop idea from the previous section in one important
way: a feedback loop can, in principle, be broken by changing what data gets
logged or by holding the model fixed for a while, since the loop runs through retraining
on new data. A performative effect does not need a retraining step to bite; it operates
the moment the model is deployed and its output starts influencing behavior,
which means even a model that is never retrained again can still be actively
distorting the outcome it was built to describe, simply by virtue of being in production
and observed by the agents it scores. Once a prediction is performative, $\Pr(Y \mid
X)$ is no longer a fixed target that a model can be fit to and then left alone;
it becomes a function of which model is deployed, since deploying a different model
changes the behavior that generates $Y$ in the first place, a possibility that
standard supervised learning, built around the assumption of a fixed target distribution,
has no native way to represent

% From: book_springer/lectures_source/Lesson02.1_From_Data_Science_To_Decision_Science.smd:624 '## Missing Exploration and Long Horizons'
\subsection{Missing Exploration and Long Horizons}
\label{sec:missing-exploration-long-horizons}

The final three sections of this chapter turn to what a model cannot know
because it was never allowed to find out: the effect of actions nobody has tried,
the consequences that only appear long after the decision was made, and the responses
of people who learn how the model scores them and adjust their behavior accordingly

% From: book_springer/lectures_source/Lesson02.1_From_Data_Science_To_Decision_Science.smd:626 '* No Exploration: Trapped in the Logged Policy'
\textbf{No Exploration: Trapped in the Logged Policy} --- A model trained purely
on observational data never experiments; it only ever sees the outcomes of
whatever policy was actually in force when the training data was logged, and it stays
trapped inside the support of that logged policy no matter how much data accumulates
It cannot learn the effect of an action nobody has ever taken, a gap known as an
overlap violation, or a positivity violation, in the causal inference literature

A pricing model illustrates the limitation cleanly. If historical prices for a product
never dropped below \$20, a model trained on that history, no matter how much of
it there is, cannot say with any evidence what would happen to demand at \$15, because
the training data simply contains no observations at that price. The model can
extrapolate a curve fit to the \$20-and-above region, but that extrapolation is
a modeling assumption, not evidence, and the true demand curve below \$20 could easily
behave in a way the extrapolation does not anticipate, for instance if a lower
price triggers a qualitatively different customer segment or a different
perception of the product

The overlap violation at the heart of this example is a property of the data,
not of the model fit to it, which means no amount of architectural sophistication
closes the gap. A more flexible model fit to the same \$20-and-above history produces
a more flexible-looking extrapolation below \$20, not a better-evidenced one,
since flexibility only changes the shape of the guess, not whether the guess is grounded
in any observed outcome. The only way to generate genuine evidence about the
untried region is to enter it, which is exactly what a deliberate pricing experiment
does and what a purely observational pipeline, by construction, never does on its
own. Some business questions cannot be answered by collecting more historical
data of the same kind; they can only be answered by a deliberate intervention,
actually trying the untried price, action, or policy, on at least a controlled subset
of cases, which is exactly the kind of exploration a purely observational
pipeline never performs

% From: book_springer/lectures_source/Lesson02.1_From_Data_Science_To_Decision_Science.smd:641 '* Delayed and Long-Horizon Effects'
\textbf{Delayed and Long-Horizon Effects} --- A model optimized for an immediate
proxy metric has no visibility into downstream outcomes that only materialize
later, and consequently no way to assign credit or blame across time for the outcomes
it actually cares about. Optimizing what is easy to measure right now can
actively work against optimizing what only becomes measurable much later

A customer support team optimized for fast ticket closure demonstrates the pattern
Closing tickets quickly is easy to measure in real time, and a team optimized
against that metric closes tickets quickly, regardless of whether the underlying
issue that prompted the ticket was actually resolved. The true cost of closing
tickets without resolving the issue, repeat complaints from the same customers and
a gradual erosion of trust in the support function, shows up only months later, and
it shows up in a completely different metric, such as churn or renewal rate, than
the one the team was originally optimized against

A collections agency optimized to maximize dollars recovered this quarter shows the
same pattern with the sign reversed: it pushes borrowers into aggressive repayment
plans that look like a win the moment they are signed. Months later, a share of
those borrowers default on the plan they were pushed into, so this quarter's extra
recovery is repaid, later, with a lower lifetime repayment rate and reputational
damage that never shows up in the metric the agency was optimized against

The underlying difficulty is a credit-assignment problem: even once the long-horizon
outcome is finally observed, it is no longer obvious which of the many earlier
actions, this ticket closure, that pricing change, an unrelated product update,
deserves credit or blame for it. Immediate proxies are attractive precisely because
they sidestep this problem entirely, reporting a number the instant an action is
taken rather than waiting, potentially for months, to find out whether the
action actually helped. That immediacy is also what makes immediate proxies dangerous
as the sole basis for optimization, since a system tuned only against what can
be measured today has no way to notice that today's convenient metric and
tomorrow's real outcome have quietly started pointing in opposite directions. Short-term
proxies and long-term business outcomes can point in opposite directions, and
only one of them, the long-term outcome, is what the business actually cares
about in the end. This means a decision system needs some mechanism for crediting
today's action with tomorrow's consequence, not just today's easily observed
result

% From: book_springer/lectures_source/Lesson02.1_From_Data_Science_To_Decision_Science.smd:662 '* Strategic and Adversarial Response'
\textbf{Strategic and Adversarial Response} --- Once a model is deployed and its
scores start to matter, the people being scored have an incentive to learn how the
model works and to adjust their behavior to score better under it, whether or
not that adjustment reflects any genuine change in the underlying quality the
model was originally trying to measure. Standard machine learning assumes passive
agents whose behavior does not respond to the model; in essentially any business
setting, the agents being scored are not passive

Publishing details about how a credit-scoring model works invites applicants to
game the specific features the model relies on rather than to improve the
underlying creditworthiness the model is meant to assess. Publishing details of a
search-ranking algorithm invites keyword-stuffing and other manipulations aimed at
the ranking signal itself, rather than at improving the underlying content
quality the ranking was designed to surface. A tax-audit model that flags returns
reporting income above a threshold produces the same response once the threshold
becomes known: self-employed filers bunch their reported income just below it,
adjusting the number they report rather than the underlying income the rule was
meant to verify

This is a close relative of Goodhart's law introduced earlier in this lesson,
but it is worth separating as its own failure mode because the mechanism is
different: Goodhart's law describes what happens when an organization optimizes a
proxy directly, while strategic response describes what happens when the people
being scored optimize the proxy on their own initiative, in response to a model whose
scoring logic they have inferred or been told. The model's training data
reflects behavior recorded before the agents had this incentive or this
knowledge, so the statistical relationship the model learned between features and
outcomes is a relationship that held under a specific, and temporary, set of behavioral
conditions

The more consequential and the more transparent a deployed model is, the faster this
adaptation tends to happen, since agents then have both a strong incentive to
respond and enough information about the scoring rule to know how to respond
effectively. A model trained on how agents behaved before it was deployed can become
systematically wrong the moment those same agents learn how the model scores them,
because the training data reflects a world in which the agents had not yet
adapted, and adaptation is precisely what a deployed, consequential model
invites

% From: book_springer/lectures_source/Lesson02.1_From_Data_Science_To_Decision_Science.smd:680 '# Why This Matters'
\section{Why This Matters}
\label{chap:why-this-matters}

% From: book_springer/lectures_source/Lesson02.1_From_Data_Science_To_Decision_Science.smd:682 '## Roadmap'
\subsection{Roadmap}
\label{sec:why-this-matters-roadmap}

The lesson closes by pulling back from the four individual costs to ask why they
matter together, what they share as a root cause, and where the rest of the book
goes to address them

% From: book_springer/lectures_source/Lesson02.1_From_Data_Science_To_Decision_Science.smd:684 '* Why AI Projects Fail'
\textbf{Why AI Projects Fail} --- Most AI project failures are not, on close
inspection, modeling failures; they are framing failures, cases where the model itself
works as intended but was pointed at the wrong problem from the start. Four
recurring patterns account for the majority of these framing failures

\begin{itemize}
  \item \textbf{Misalignment}: the model optimizes a proxy metric rather than the
    actual business goal, so improvements in the trained metric do not translate
    into the outcome the business cares about

  \item \textbf{Black boxes}: the model's outputs cannot be explained in terms a
    stakeholder can evaluate, so they are not trusted and, in turn, are not
    adopted into real decision-making

  \item \textbf{No lever}: the pipeline delivers a score where the business needed
    an action, leaving the translation from score to decision unresolved

  \item \textbf{No feedback}: performance is never monitored after deployment,
    so drift, feedback loops, and strategic responses go undetected until they
    have already done damage
\end{itemize}

Each of these four patterns can coexist with a model that is, by every standard machine
learning metric, excellent, which is exactly what makes them hard to catch in a code
review or a model evaluation. Misalignment survives a rigorous train/test split,
since the proxy metric is being measured correctly, just measuring the wrong
thing. A black box can pass every accuracy benchmark while still failing the moment
a regulator, a clinician, or a customer asks for the reasoning behind a specific
decision. A missing lever is invisible in an offline evaluation, since offline evaluation
only ever checks the score against a held-out label, never against whether
anyone downstream knew what to do with the score. And the absence of feedback is,
almost by definition, undetectable at launch, since it is a gap in what happens
after launch, not in anything that can be checked before it

Recognizing these four patterns as framing failures rather than modeling
failures changes where the fix belongs: the remedy is not a better model architecture
or more training data, but a different specification of the problem before any
model is built, one that names the true objective, requires an explainable
structure, identifies the lever a stakeholder needs, and builds in monitoring
from the start. An accurate model that answers the wrong question is a failed project,
however good its metrics look on a dashboard, because the dashboard is measuring
fit to a proxy, not success at the decision the business actually needed made

% From: book_springer/lectures_source/Lesson02.1_From_Data_Science_To_Decision_Science.smd:698 '* The Root Cause of the Four Costs'
\textbf{The Root Cause of the Four Costs} --- The four costs surveyed across this
lesson, ignoring causality, ignoring uncertainty, ignoring the business objective,
and ignoring dynamics, can look like four unrelated engineering gaps, each
requiring its own separate fix. They share a single root cause. Each one arises
from optimizing prediction accuracy on a fixed distribution, instead of optimizing
decision quality on a distribution that the decision itself changes once it is acted
upon. This book addresses all four costs together with the same causal and
probabilistic reasoning: every technique introduced in a later chapter plugs into
one recurring loop,
\begin{equation}
  \label{eq:decision-loop}\text{data}\to \text{model}\to \text{policy}\to \text{action}
  \to \text{feedback}\;,
\end{equation}
with the feedback arrow closing back on the data that trains the next model

Causality is ignored because a fixed-distribution predictor never needs to
distinguish correlation from mechanism to fit its training data well, as Chapter~\ref{chap:cost-ignoring-causality}
showed with confounding, missing counterfactuals, and the structural failure modes
that survive even after confounding is addressed. Uncertainty is ignored because
a single point estimate is sufficient to minimize average error on a fixed
distribution, even though it says nothing about how that estimate would hold up under
a different one, as Chapter~\ref{chap:cost-ignoring-uncertainty} showed with
error bars, the aleatoric-epistemic split, and the absence of abstention. The business
objective is ignored because any proxy that correlates with the true goal on the
fixed training distribution looks like a reasonable target to optimize, as
Chapter~\ref{chap:cost-ignoring-business-objective} showed with Goodhart's law, asymmetric
losses, and the cost of splitting a decision into separately optimized pieces. Dynamics
are ignored because a fixed-distribution assumption is, by definition, an assumption
that the distribution does not move, which is precisely the assumption a
deployed decision violates, as Chapter~\ref{chap:cost-ignoring-dynamics-feedback}
showed with distribution shift, feedback loops, and strategic response

Seen this way, the four chapters are not four independent lessons to be learned separately;
they are four consequences of the same modeling choice, examined from four different
angles. Fixing all four costs at once is not four separate engineering projects
bolted onto an existing pipeline; it is a single shift in what the pipeline is built
to optimize, from prediction accuracy on data already collected to decision
quality on the consequences the decision itself sets in motion

% From: book_springer/lectures_source/Lesson02.1_From_Data_Science_To_Decision_Science.smd:754 '* Book Roadmap'
\textbf{Book Roadmap} --- This lesson has catalogued what a purely predictive
pipeline misses: causal mechanism, calibrated uncertainty, an explicit business
objective, and an accounting for the dynamics a decision sets off. The rest of
the book builds the pipeline that does not miss these things, organized into
four parts that build on one another

\begin{itemize}
  \item \textbf{Part~II: Advanced Modeling Theory \& Tools} supplies the
    machinery the later parts depend on: knowledge representation, probabilistic
    machine learning (Bayesian inference, uncertainty), and causal machine
    learning (structural causal models, do-calculus, causal discovery)

  \item \textbf{Part~III: Data} turns domain knowledge into causal DAGs and
    causal data pipelines, and confronts the selection bias and distribution
    shift surveyed in this chapter

  \item \textbf{Part~IV: Decision-Making Theory \& Tools} turns causal models
    into actions: decision theory, the taxonomy of decision problems, simple
    decisions such as bandits and reinforcement learning, complex decisions such
    as policy gradients and multi-agent settings, and agentic causal reasoning

  \item \textbf{Part~V: Implementation, Deployment \& Governance} covers
    stakeholder alignment, deployment, monitoring, and adaptation, and trust,
    explainability, fairness, and governance
\end{itemize}

Each part picks up where the previous one leaves off. Part~II's probabilistic and
causal tools supply the representational machinery, a way to express uncertainty
and a way to express causal structure, that the earlier chapters showed a standard
predictive pipeline lacks. Part~III then turns that machinery on the data itself,
building the causal DAGs and data pipelines a decision needs while confronting
head-on the selection bias and distribution shift that Chapter~\ref{chap:cost-ignoring-causality}
and Chapter~\ref{chap:cost-ignoring-dynamics-feedback} showed a purely observational
pipeline cannot see past on its own. Part~IV then puts the machinery and the data
to work choosing actions, from a single simple decision to the multi-step, multi-agent,
and agentic settings that the feedback loops and strategic responses of
Chapter~\ref{chap:cost-ignoring-dynamics-feedback} demand. Part~V closes the loop
back to practice, since a decision system that is well designed on paper still has
to be deployed, watched, adjusted, and governed once it is running against real
data it can no longer fully control, and once it is answering to the stakeholders
Chapter~\ref{chap:cost-ignoring-business-objective} showed a bare score cannot
satisfy on its own

Every later technique in this book is best understood as strengthening one arrow
in the loop of Eq.~\eqref{eq:decision-loop}, and the loop closing back on itself,
through the feedback arrow, is exactly the dynamic that Sect.~\ref{sec:world-that-reacts}
showed a purely predictive pipeline has no way to represent
